\documentclass[11pt]{article}

\usepackage[T1]{fontenc}
\usepackage{lmodern}
\usepackage{amsmath,amssymb,amsthm,mathtools}
\usepackage[margin=1in]{geometry}
\usepackage{xurl}
\usepackage[hidelinks]{hyperref}

\newtheorem{theorem}{Theorem}[section]
\newtheorem{proposition}[theorem]{Proposition}
\newtheorem{corollary}[theorem]{Corollary}
\theoremstyle{remark}
\newtheorem{remark}[theorem]{Remark}

\newcommand{\Ncurv}{\mathcal{N}_{\varepsilon}}
\newcommand{\NN}[1]{\mathcal{N}_{\varepsilon,#1}}
\newcommand{\R}{\mathbb{R}}
\newcommand{\C}{\mathbb{C}}
\newcommand{\Om}{\Omega}

\title{M6B: Nonuniqueness of the Regularized Nonlinear Curvature Closure\\
\large Two Exact Counterexamples and Their Scope}
\author{Richard J. Reyes}
\date{October 26, 2026}

\begin{document}
\maketitle

\begin{abstract}
We examine the proposed nonlinear curvature operator
\[
\mathcal{N}_{\varepsilon}[\psi]
=-\frac{(\Delta\psi)\overline{\psi}}
       {|\psi|^2+\varepsilon^2e^{-2\alpha|\psi|^2}},
\qquad \varepsilon>0,\quad\alpha\in\R.
\]
Two distinct uniqueness statements fail. First, positivity, smoothness at
nodes, invariance under constant phase rotations, and recovery of the
unregularized curvature quotient do not characterize the closure uniquely:
we construct an infinite family of different operators with all these
properties. Second, on a bounded square with homogeneous Dirichlet data,
we construct two smooth complex fields having the same prescribed positive
mass and the same nonlinear curvature image, although they are not related
by any constant phase. These are counterexamples to uniqueness of the
closure and injectivity of the stationary curvature map, respectively.
They do not imply nonuniqueness of solutions of the time-dependent
curvature--wavefield initial-value problem.
\end{abstract}

\section{Definition and preliminary facts}

Let $d\geq1$, let $\Omega\subset\R^d$ be open, and let
$\psi\in C^2(\Omega;\C)$. For $\varepsilon>0$ and $\alpha\in\R$, set
\begin{equation}\label{eq:denominator}
D_\varepsilon(s)=s+\varepsilon^2e^{-2\alpha s},
\qquad s\geq0,
\end{equation}
and define the canonical M6B operator by
\begin{equation}\label{eq:canonical}
\boxed{\quad
\Ncurv[\psi]
=-\frac{(\Delta\psi)\overline\psi}
       {D_\varepsilon(|\psi|^2)}.
\quad}
\end{equation}

\begin{proposition}[Node safety and the unregularized limit]
\label{prop:basic}
For every $s\geq0$, $D_\varepsilon(s)>0$. Consequently,
\eqref{eq:canonical} is pointwise defined for every $C^2$ field, including
at zeros of $\psi$, where $\Ncurv[\psi]=0$. At each point where
$\psi\neq0$,
\begin{equation}\label{eq:unreg}
\lim_{\varepsilon\downarrow0}\Ncurv[\psi]
=-\frac{\Delta\psi}{\psi}.
\end{equation}
For every constant $\theta\in\R$,
\begin{equation}\label{eq:invariance}
\Ncurv[e^{i\theta}\psi]=\Ncurv[\psi].
\end{equation}
\end{proposition}

\begin{proof}
Since $s\geq0$, $\varepsilon^2>0$, and $e^{-2\alpha s}>0$,
one has $D_\varepsilon(s)>0$. At $\psi=0$, the numerator in
\eqref{eq:canonical} vanishes while the denominator equals
$\varepsilon^2$. At a point with $\psi\neq0$, the denominator converges
to $|\psi|^2$ as $\varepsilon\downarrow0$, and
$\overline\psi/|\psi|^2=1/\psi$, giving \eqref{eq:unreg}.
Finally, constancy of $\theta$ implies
$\Delta(e^{i\theta}\psi)=e^{i\theta}\Delta\psi$; the phases cancel in
the numerator, and the modulus in the denominator is unchanged.
\end{proof}

\section{Nonuniqueness of the regularized closure}

\begin{theorem}[The stated structural axioms do not select a unique closure]
\label{thm:closure}
Fix $d\geq1$, $\varepsilon>0$, $\alpha\in\R$, and a reference
intensity $S_0>0$. There exists an infinite family of pairwise distinct
operators indexed by $\lambda\geq0$ such that each operator:
\begin{enumerate}
\item has a strictly positive denominator and is smooth in the real
      field variables and their second derivatives;
\item is well-defined and vanishes at field nodes;
\item is invariant under constant global $U(1)$ phase rotations;
\item converges pointwise to $-\Delta\psi/\psi$ as
      $\varepsilon\downarrow0$ wherever $\psi\neq0$;
\item has denominator $\varepsilon^2$ at $\psi=0$.
\end{enumerate}
The family includes the canonical M6B operator.
\end{theorem}

\begin{proof}
For $s\geq0$, define
\begin{equation}\label{eq:family-denominator}
D_{\varepsilon,\lambda}(s)
=s+\varepsilon^2e^{-2\alpha s}
  \left(1+\lambda\frac{s}{S_0+s}\right),
\qquad \lambda\geq0,
\end{equation}
and for $\psi\in C^2(\Omega;\C)$ put
\begin{equation}\label{eq:family}
\NN{\lambda}[\psi]
=-\frac{(\Delta\psi)\overline\psi}
       {D_{\varepsilon,\lambda}(|\psi|^2)}.
\end{equation}
When $\lambda=0$, \eqref{eq:family} equals \eqref{eq:canonical}.

\emph{Positivity and smoothness.}
For $s\geq0$ and $\lambda\geq0$, the multiplier
$1+\lambda s/(S_0+s)$ is at least one. Hence
\[
D_{\varepsilon,\lambda}(s)
\geq \varepsilon^2e^{-2\alpha s}>0.
\]
All functions in \eqref{eq:family-denominator} are smooth on an open
neighborhood of the half-line $s\geq0$, and the denominator is nonzero
there. Therefore the jet-variable map
\[
(z,w)\longmapsto
-\frac{w\overline z}{D_{\varepsilon,\lambda}(|z|^2)}
\qquad (z,w\in\C)
\]
is smooth. Substitution of $z=\psi(x)$ and $w=\Delta\psi(x)$ gives
pointwise smoothness of the closure.

\emph{Nodes.}
At $s=0$,
\[
D_{\varepsilon,\lambda}(0)=\varepsilon^2.
\]
When $\psi(x)=0$, its curvature image at $x$ is zero, regardless of
$\Delta\psi(x)$.

\emph{Global phase invariance.}
For any constant $\theta\in\R$,
\[
\Delta(e^{i\theta}\psi)\,
\overline{e^{i\theta}\psi}
=e^{i\theta}\Delta\psi\,e^{-i\theta}\overline\psi
=(\Delta\psi)\overline\psi,
\]
and $|e^{i\theta}\psi|^2=|\psi|^2$. Equation
\eqref{eq:family} is consequently unchanged.

\emph{Unregularized limit.}
At any point where $\psi\neq0$, put $s=|\psi|^2>0$. For each fixed
$\lambda\geq0$,
\[
\lim_{\varepsilon\downarrow0}D_{\varepsilon,\lambda}(s)=s.
\]
It follows that
\[
\lim_{\varepsilon\downarrow0}\NN{\lambda}[\psi]
=-\frac{(\Delta\psi)\overline\psi}{|\psi|^2}
=-\frac{\Delta\psi}{\psi}.
\]

\emph{Pairwise distinctness.}
For $s>0$ and $\lambda_1\neq\lambda_2$,
\begin{equation}\label{eq:difference}
D_{\varepsilon,\lambda_1}(s)
-D_{\varepsilon,\lambda_2}(s)
=(\lambda_1-\lambda_2)\varepsilon^2e^{-2\alpha s}
  \frac{s}{S_0+s}\neq0.
\end{equation}
Choose a neighborhood of $0\in\R^d$ and the smooth test field
$\psi(x)=1+|x|^2$. Then
$\psi(0)=1$ and $\Delta\psi(0)=2d$. Hence
\begin{equation}\label{eq:distinct}
\NN{\lambda}[\psi](0)
=-\frac{2d}
{1+\varepsilon^2e^{-2\alpha}
 \left(1+\lambda/(S_0+1)\right)}.
\end{equation}
The right-hand side takes distinct values at distinct $\lambda\geq0$.
Therefore all operators in the family are different.
\end{proof}

\begin{remark}
Theorem~\ref{thm:closure} refutes uniqueness \emph{under the enumerated
properties}. It does not rule out a uniqueness result derived from an
additional, explicitly specified action principle, normalization,
constitutive law, or physical constraint.
\end{remark}

\section{Noninjectivity of the canonical curvature map modulo phase}

\begin{theorem}[Two inequivalent fields with identical curvature,
Dirichlet data, and mass]
\label{thm:stationary}
Let $\Omega=(0,\pi)^2$, $\varepsilon>0$, and $\alpha\in\R$.
For every $m>0$, there exists a real-valued
$F_m\in C^\infty(\overline\Omega)$ such that
\begin{equation}\label{eq:stationary-pde}
\begin{cases}
\Ncurv[\psi]=F_m&\text{in }\Omega,\\
\psi=0&\text{on }\partial\Omega,\\
\|\psi\|_{L^2(\Omega)}^2=m
\end{cases}
\end{equation}
has at least two solutions
$\psi_\pm\in C^\infty(\overline\Omega;\C)$ with
\[
\psi_-\neq e^{i\theta}\psi_+
\quad\text{for every constant }\theta\in\R.
\]
Thus the canonical stationary curvature map is not injective even
modulo global phase, with mass and homogeneous Dirichlet data fixed.
\end{theorem}

\begin{proof}
\emph{Common Laplacian eigenvalue.}
Let
\[
f(x,y)=\sin x\sin(2y),\qquad
 g(x,y)=\sin(2x)\sin y.
\]
Both are real and smooth on $\overline\Omega$ and vanish on
$\partial\Omega$. Direct differentiation yields
\begin{equation}\label{eq:eigen}
\Delta f=-5f,\qquad\Delta g=-5g.
\end{equation}

\emph{Two complex fields.}
For arbitrary prescribed $m>0$, set
\[
c=\frac{\sqrt{2m}}{\pi}>0,
\qquad
\psi_+=c(f+ig),\quad
\psi_-=c(f-ig).
\]
By \eqref{eq:eigen},
\begin{equation}\label{eq:eigen-fields}
\Delta\psi_\pm=-5\psi_\pm,
\qquad\psi_\pm|_{\partial\Omega}=0.
\end{equation}
Moreover, since $f$ and $g$ are real,
\begin{equation}\label{eq:amplitude}
|\psi_+|^2=|\psi_-|^2
=c^2(f^2+g^2)=:A_m(x,y).
\end{equation}

\emph{Identical curvature images.}
Using \eqref{eq:eigen-fields} in the canonical formula,
\begin{align*}
\Ncurv[\psi_\pm]
&=-\frac{(-5\psi_\pm)\overline{\psi_\pm}}
{|\psi_\pm|^2+\varepsilon^2e^{-2\alpha|\psi_\pm|^2}}\\
&=\frac{5|\psi_\pm|^2}
{|\psi_\pm|^2+\varepsilon^2e^{-2\alpha|\psi_\pm|^2}}\\
&=\frac{5A_m(x,y)}{A_m(x,y)+\varepsilon^2e^{-2\alpha A_m(x,y)}}.
\end{align*}
Define
\begin{equation}\label{eq:force}
\boxed{\quad
F_m(x,y)=\frac{5A_m(x,y)}
{A_m(x,y)+\varepsilon^2e^{-2\alpha A_m(x,y)}}.
\quad}
\end{equation}
The positive denominator proves that $F_m$ is real-valued and
$C^\infty$ on $\overline\Omega$. Both fields solve the same equation
with the same source.

\emph{Equal prescribed masses.}
For every integer $n\geq1$,
\[
\int_0^\pi\sin^2(nx)\,dx=\frac\pi2.
\]
Consequently,
\[
\|f\|_2^2=\frac{\pi^2}{4},\qquad
\|g\|_2^2=\frac{\pi^2}{4},
\]
and therefore
\begin{equation}\label{eq:mass}
\|\psi_\pm\|_2^2
=c^2\int_\Omega(f^2+g^2)\,dx\,dy
=\frac{2m}{\pi^2}\frac{\pi^2}{2}=m.
\end{equation}

\emph{Not equivalent under any global phase.}
Suppose that $\psi_-=e^{i\theta}\psi_+$ for a constant $\theta$.
At
\[
P=(\pi/2,\pi/4),
\]
we have $f(P)=1$ and $g(P)=0$. Thus
$\psi_+(P)=\psi_-(P)=c\neq0$, forcing $e^{i\theta}=1$.
At
\[
Q=(\pi/4,\pi/2),
\]
we have $f(Q)=0$ and $g(Q)=1$. Hence
$\psi_+(Q)=ic$ and $\psi_-(Q)=-ic$, forcing
$e^{i\theta}=-1$. This is a contradiction.

The two fields are therefore inequivalent modulo global phase and
satisfy the same boundary conditions, source equation, and mass
constraint. This proves the theorem.
\end{proof}

\section{Interpretation and limitations}

\begin{proposition}[Phase invariance is not phase equivariance]
For any $\psi\in C^2(\Omega;\C)$ and constant $\theta\in\R$,
\[
\Ncurv[e^{i\theta}\psi]=\Ncurv[\psi].
\]
In general, however,
\[
\Ncurv[e^{i\theta}\psi]\neq
 e^{i\theta}\Ncurv[\psi].
\]
\end{proposition}
\begin{proof}
The first identity is Proposition~\ref{prop:basic}. For the second,
take any field with $\Ncurv[\psi](x)\neq0$ at some point $x$, for
example $\psi(x)=1+|x|^2$ near $0$. If $e^{i\theta}\neq1$, the two
expressions at that point differ.
\end{proof}

\begin{remark}[Precise meaning of the counterexamples]
Theorem~\ref{thm:closure} shows that the stipulated regularity,
symmetry, and limit conditions do not identify one \emph{constitutive
operator}. Theorem~\ref{thm:stationary} shows that the fixed canonical
mapping $\psi\mapsto\Ncurv[\psi]$ is not injective on a particular
Dirichlet and fixed-mass class. Neither theorem states that a
specified time-dependent PDE has two solutions with the \emph{same
initial condition}. Such a conclusion would require a separate
analysis of the chosen evolution equation, its domain, regularity
class, parameters, and initial data. Similarly, existence for an
arbitrary prescribed forcing $F$ is not established here.
\end{remark}

\section{Conclusion}

The correct classification for the claims proved here is
\[
\boxed{\text{M6B: explicit counterexamples to closure uniqueness
and stationary injectivity.}}
\]
Both counterexamples are analytical, require no numerical search, and
hold for all $\varepsilon>0$ and $\alpha\in\R$ under their stated
hypotheses. Existence and uniqueness for a separately specified full
nonlinear evolution remain distinct open obligations.

\begin{thebibliography}{9}
\bibitem{registry}
\emph{WCT Full Equation List (Corrected)}, M6B: nonlinear curvature
operator. \href{https://github.com/rickyjreyes/geometry_of_resonance/blob/main/WCT_FULL_EQUATION_LIST_CORRECTED.md}{Canonical registry document}
\bibitem{issue}
\emph{wct-sympy}, issue \#17: M6B nonlinear curvature-operator
existence/uniqueness obligations.
\href{https://github.com/rickyjreyes/wct-sympy/issues/17}{M6B issue tracker}
\end{thebibliography}

\end{document}
